\section{Convenzioni per il testo}

\paragraph{Numeri}
Invece di scriverli, i numeri venivano scritti a lettere.

\paragraph{Spazi}
Per evitare facilitazioni di analisi tramite la lunghezza delle parole,
non venivano inseriti spazi.
Al loro posto, lo spazio veniva rimpiazato da una "X".

\paragraph{Fine paragrafo}
Similmente agli spazi, veniva usato una "XX".

\paragraph{Virgola}
Veniva usata una "Y".

\paragraph{Trattino}
Veniva usata una "YY".

\paragraph{Parentesi}
Veniva usata una "KK" sia per aprire che per chiudere.

\paragraph{Apici}
Veniva usata una "J".

\paragraph{Punto di domanda}
Veniva usata una "UD".

\paragraph{Sostituzione coppie di lettere}
Era possibile, nei sottomarini, invece di scrivere il testo in chiaro, cifrare direttamente una coppia di lettere.
Questo era possibile tramite una tabella che consentiva di sostiturire una coppia di lettere con una seconda.

\paragraph{Nomi importanti}
Invece di scrivere per esteso i nomi importanti, essi venivano
abbreviati.

\paragraph{Ridondanza testo}
Per evitare ricezione errate dei messaggi, il quartier generarle (marina), venivano ripetuti.
4 volte in un ora dal primo messaggio. Uno ogni 15'.

\paragraph{Limite caratteri}
Per evitare messaggi troppo lunghi 250 caratteri fu posto come limite per i messaggi.
